\chapter{Chapter 5: Vector Calculus}

\begin{enumerate}

%%% PROBLEM 1 %%%
\item \textbf{Solution.}
\begin{enumerate}
\item We substitute $x=1$, $y=-1$:
\[
\mathbf{F}(1,-1) = (1+1)\,\mathbf{i} + (-1-2)\,\mathbf{j} = 2\,\mathbf{i} - 3\,\mathbf{j}.
\]

\item At $(0,0)$: $\mathbf{F}=(1,-2)$. At $(1,0)$: $\mathbf{F}=(2,-2)$. At $(0,1)$: $\mathbf{F}=(1,-1)$. At $(-1,0)$: $\mathbf{F}=(0,-2)$. At $(0,-1)$: $\mathbf{F}=(1,-3)$.

The field has a ``center'' at $(-1,2)$ where $\mathbf{F}=\mathbf{0}$. Near the origin, every vector points roughly to the lower-right, with the $y$-component always negative (pulling downward) and the $x$-component positive for $x>-1$. The field radiates outward from the point $(-1,2)$.
\end{enumerate}

%%% PROBLEM 2 %%%
\item \textbf{Solution.}
A \emph{scalar field} assigns a single real number to each point in space. Example: the temperature $T(x,y,z)$ at each point in a room.

A \emph{vector field} assigns a vector to each point in space. Example: the velocity $\mathbf{v}(x,y,z)$ of a fluid at each point in a flowing stream.

The key difference is that a scalar field has only magnitude at each point, while a vector field has both magnitude and direction.

%%% PROBLEM 3 %%%
\item \textbf{Solution.}
At $(2,0)$: $\mathbf{F} = -0\,\mathbf{i} + 2\,\mathbf{j} = (0,2)$, so the arrow points straight up.

At $(1,1)$: $\mathbf{F} = -1\,\mathbf{i} + 1\,\mathbf{j} = (-1,1)$, so the arrow points up and to the left at $45^\circ$.

At $(-2,-2)$: $\mathbf{F} = 2\,\mathbf{i} + (-2)\,\mathbf{j} = (2,-2)$, so the arrow points down and to the right.

In each case, $\mathbf{F}$ is perpendicular to the position vector $(x,y)$ since $\mathbf{F}\cdot(x,y)= -yx + xy = 0$. Also $|\mathbf{F}| = \sqrt{x^2+y^2}$, which equals the distance from the origin. The overall pattern is \emph{counterclockwise rotation} about the origin, with speed proportional to distance.

%%% PROBLEM 4 %%%
\item \textbf{Solution.}
\begin{enumerate}
\item $\mathbf{F}(0,0) = (e^0, e^0) = (1,1)$.

$\mathbf{F}(1,2) = (e^1, e^2) = (e, e^2) \approx (2.718, 7.389)$.

\item At the origin, $\mathbf{F}(0,0)=(1,1)$ is nonzero and points into the first quadrant. Since both components $e^x > 0$ and $e^y > 0$ everywhere, every vector points into the first quadrant. The origin is \emph{neither} a source nor a sink, because $\mathbf{F}(0,0)\neq\mathbf{0}$; sources and sinks occur at points where $\mathbf{F}=\mathbf{0}$ and the flow radiates outward or inward, respectively.
\end{enumerate}

%%% PROBLEM 5 %%%
\item \textbf{Solution.}
Each component of $\mathbf{G}$ has $x^2+y^2+z^2$ in the denominator, which is zero only at the origin $(0,0,0)$. For any point $(x,y,z)\neq(0,0,0)$, the components are ratios of polynomials with a nonzero denominator, hence continuous.

Therefore, $\mathbf{G}$ is \emph{not} continuous at the origin $(0,0,0)$ (it is undefined there), but it is continuous everywhere else in $\mathbb{R}^3\setminus\{(0,0,0)\}$.

%%% PROBLEM 6 %%%
\item \textbf{Solution.}
\begin{enumerate}
\item Each component $x^2$, $y^2$, $z^2$ is a polynomial, and polynomials are continuous on all of $\mathbb{R}^3$. Therefore $\mathbf{H}$ is continuous at every point in $\mathbb{R}^3$.

\item There are no problematic points. Every component is a polynomial, which is defined and continuous everywhere.
\end{enumerate}

%%% PROBLEM 7 %%%
\item \textbf{Solution.}
Engineers and physicists sometimes normalize all arrows to the same length to make the \emph{direction} of the field clearly visible at every point, especially in regions where the magnitude varies dramatically.

\textbf{Advantage:} Normalized plots clearly show the directional pattern of the field (e.g., convergence, divergence, rotation) without large vectors obscuring small ones.

\textbf{Disadvantage:} All information about the \emph{magnitude} is lost. For example, in a wind field, one cannot distinguish regions of high wind speed from regions of low wind speed.

%%% PROBLEM 8 %%%
\item \textbf{Solution.}
\begin{enumerate}
\item $\mathbf{F}_1(x,y) = x\,\mathbf{i} + y\,\mathbf{j}$: vectors point radially \emph{outward} from the origin. This is a \textbf{source}.

\item $\mathbf{F}_2(x,y) = -x\,\mathbf{i} - y\,\mathbf{j}$: vectors point radially \emph{inward} toward the origin. This is a \textbf{sink}.

\item $\mathbf{F}_3(x,y) = x\,\mathbf{i} - y\,\mathbf{j}$: vectors point outward along the $x$-axis but inward along the $y$-axis (or vice versa). This is a \textbf{saddle}.
\end{enumerate}

%%% PROBLEM 9 %%%
\item \textbf{Solution.}
Write $\mathbf{F}(x,y) = P\,\mathbf{i} + Q\,\mathbf{j}$ where $P = y$ and $Q = -x$.

Compute:
\[
\frac{\partial P}{\partial y} = 1, \qquad \frac{\partial Q}{\partial x} = -1.
\]
Since $\dfrac{\partial P}{\partial y} = 1 \neq -1 = \dfrac{\partial Q}{\partial x}$, the field is \textbf{not conservative}.

The failure of the conservative test corresponds to the fact that $\mathbf{F}$ has a rotational/circulatory pattern near the origin---vectors circulate (clockwise) around the origin rather than emanating from a potential function.

%%% PROBLEM 10 %%%
\item \textbf{Solution.}
\begin{enumerate}
\item Here $P = x^3$ and $Q = 3x^2y$.
\[
\frac{\partial P}{\partial y} = 0, \qquad \frac{\partial Q}{\partial x} = 6xy.
\]
Since $0 \neq 6xy$ in general, $\mathbf{F}$ is \textbf{not conservative}.

\item Since $\mathbf{F}$ is not conservative, no potential function exists.
\end{enumerate}

%%% PROBLEM 11 %%%
\item \textbf{Solution.}
Write $\mathbf{F} = (P,Q,R) = (2xy,\; x^2+z^2,\; 2yz)$. For $\mathbf{F}$ to be conservative on $\mathbb{R}^3$ we need:
\[
\frac{\partial P}{\partial y} = \frac{\partial Q}{\partial x}, \quad \frac{\partial P}{\partial z} = \frac{\partial R}{\partial x}, \quad \frac{\partial Q}{\partial z} = \frac{\partial R}{\partial y}.
\]
Compute each:
\[
\frac{\partial P}{\partial y} = 2x, \quad \frac{\partial Q}{\partial x} = 2x.
\]
\[
\frac{\partial P}{\partial z} = 0, \quad \frac{\partial R}{\partial x} = 0.
\]
\[
\frac{\partial Q}{\partial z} = 2z, \quad \frac{\partial R}{\partial y} = 2z.
\]
All three conditions are satisfied, so $\mathbf{F}$ \textbf{is conservative} on $\mathbb{R}^3$.

To find the potential function: $f_x = 2xy \Rightarrow f = x^2y + g(y,z)$. Then $f_y = x^2 + g_y = x^2 + z^2$, so $g_y = z^2$, giving $g = yz^2 + h(z)$. Then $f_z = 2yz + h'(z) = 2yz$, so $h'(z)=0$ and $h$ is constant. Thus $f(x,y,z) = x^2y + yz^2 + C$.

%%% PROBLEM 12 %%%
\item \textbf{Solution.}
We have $P = \sqrt{x^2+y^2}$ and $Q = -\sin(xy)$.

\[
\frac{\partial P}{\partial x} = \frac{x}{\sqrt{x^2+y^2}}, \qquad
\frac{\partial P}{\partial y} = \frac{y}{\sqrt{x^2+y^2}}.
\]
\[
\frac{\partial Q}{\partial x} = -y\cos(xy), \qquad
\frac{\partial Q}{\partial y} = -x\cos(xy).
\]

At the origin $(0,0)$: $\sqrt{x^2+y^2}=0$, so $\dfrac{\partial P}{\partial x}$ and $\dfrac{\partial P}{\partial y}$ involve division by zero. These two partial derivatives are \textbf{undefined at the origin}.

The partial derivatives $\dfrac{\partial Q}{\partial x} = -y\cos(xy)$ and $\dfrac{\partial Q}{\partial y} = -x\cos(xy)$ are well-defined everywhere, including at the origin (both equal $0$ there).

%%% PROBLEM 13 %%%
\item \textbf{Solution.}
\begin{enumerate}
\item $P = x^2+y^2$ and $Q = -(x^2+y^2)$.
\[
\frac{\partial P}{\partial y} = 2y, \qquad \frac{\partial Q}{\partial x} = -2x.
\]
Since $2y \neq -2x$ in general, $\mathbf{F}$ is \textbf{not conservative}.

\item Since the field is not conservative, we classify the origin by examining the flow. At the origin $\mathbf{F}(0,0)=(0,0)$, so it is an equilibrium point. Near the origin, $\text{div}(\mathbf{F}) = \frac{\partial P}{\partial x} + \frac{\partial Q}{\partial y} = 2x + (-2y) = 2x - 2y$. At the origin, $\text{div}(\mathbf{F})=0$, and the divergence changes sign depending on direction. Examining nearby vectors: along the positive $x$-axis, $\mathbf{F}=(x^2, -x^2)$ points into the fourth quadrant; along the positive $y$-axis, $\mathbf{F}=(y^2, -y^2)$ also points into the fourth quadrant. The flow pattern is neither purely radial inward nor outward---the origin behaves like a \textbf{saddle} (flow is neither a pure source nor a pure sink).
\end{enumerate}

%%% PROBLEM 14 %%%
\item \textbf{Solution.}
Parameterize $C$: $\mathbf{r}(t) = (2t, 4t)$ for $0 \le t \le 1$.

Then $\mathbf{r}'(t) = (2,4)$ and $\|\mathbf{r}'(t)\| = \sqrt{4+16} = \sqrt{20} = 2\sqrt{5}$.

On $C$: $x+y = 2t + 4t = 6t$.
\[
\int_C (x+y)\,ds = \int_0^1 6t \cdot 2\sqrt{5}\,dt = 12\sqrt{5}\int_0^1 t\,dt = 12\sqrt{5}\cdot\frac{1}{2} = \boxed{6\sqrt{5}}.
\]

%%% PROBLEM 15 %%%
\item \textbf{Solution.}
Parameterize: $\mathbf{r}(t) = (t, t^2)$ for $0 \le t \le 2$.

$\mathbf{r}'(t) = (1, 2t)$, so $\|\mathbf{r}'(t)\| = \sqrt{1+4t^2}$.

On $C$: $x^2+y = t^2 + t^2 = 2t^2$.
\[
\int_C (x^2+y)\,ds = \int_0^2 2t^2\sqrt{1+4t^2}\,dt.
\]
Let $u = 1+4t^2$, then $du = 8t\,dt$. This substitution does not directly simplify $t^2\sqrt{1+4t^2}$. Instead, use the substitution $2t = \sinh\theta$, so $4t^2 = \sinh^2\theta$, $\sqrt{1+4t^2}=\cosh\theta$, and $2\,dt = \cosh\theta\,d\theta$.

When $t=0$: $\theta=0$. When $t=2$: $\sinh\theta=4$, so $\theta = \ln(4+\sqrt{17})$.

\[
\int_0^2 2t^2\sqrt{1+4t^2}\,dt = \int_0^{\theta_0} 2\cdot\frac{\sinh^2\theta}{4}\cdot\cosh\theta\cdot\frac{\cosh\theta}{2}\,d\theta = \frac{1}{4}\int_0^{\theta_0}\sinh^2\theta\cosh^2\theta\,d\theta
\]
where $\theta_0 = \sinh^{-1}(4)$. Using $\sinh^2\theta\cosh^2\theta = \frac{1}{4}\sinh^2(2\theta)= \frac{1}{8}(\cosh(4\theta)-1)$:
\[
= \frac{1}{4}\cdot\frac{1}{8}\left[\frac{\sinh(4\theta)}{4}-\theta\right]_0^{\theta_0} = \frac{1}{32}\left[\frac{\sinh(4\theta_0)}{4}-\theta_0\right].
\]
Now $\sinh\theta_0 = 4$, $\cosh\theta_0 = \sqrt{17}$. We need $\sinh(4\theta_0)$. Using double-angle formulas: $\sinh(2\theta_0)=2\cdot4\cdot\sqrt{17}=8\sqrt{17}$, $\cosh(2\theta_0)=2\cosh^2\theta_0-1=33$. Then $\sinh(4\theta_0) = 2\sinh(2\theta_0)\cosh(2\theta_0) = 2\cdot8\sqrt{17}\cdot33 = 528\sqrt{17}$.

\[
= \frac{1}{32}\left[\frac{528\sqrt{17}}{4} - \ln(4+\sqrt{17})\right] = \frac{1}{32}\left[132\sqrt{17} - \ln(4+\sqrt{17})\right].
\]

\[
\boxed{\int_C (x^2+y)\,ds = \frac{132\sqrt{17} - \ln(4+\sqrt{17})}{32}.}
\]

%%% PROBLEM 16 %%%
\item \textbf{Solution.}
Let $\mathbf{r}(t)$, $a\le t\le b$, be a parameterization of $C$. The scalar line integral is
\[
\int_C f\,ds = \int_a^b f(\mathbf{r}(t))\,\|\mathbf{r}'(t)\|\,dt.
\]
Now reverse the parameterization: let $\tilde{\mathbf{r}}(s) = \mathbf{r}(a+b-s)$ for $s\in[a,b]$. Then $\tilde{\mathbf{r}}\,'(s) = -\mathbf{r}'(a+b-s)$, so
\[
\|\tilde{\mathbf{r}}\,'(s)\| = \|\mathbf{r}'(a+b-s)\|.
\]
The integral becomes:
\[
\int_a^b f(\tilde{\mathbf{r}}(s))\,\|\tilde{\mathbf{r}}\,'(s)\|\,ds = \int_a^b f(\mathbf{r}(a+b-s))\,\|\mathbf{r}'(a+b-s)\|\,ds.
\]
Substituting $u = a+b-s$, $du = -ds$; when $s=a$, $u=b$ and when $s=b$, $u=a$:
\[
= \int_b^a f(\mathbf{r}(u))\,\|\mathbf{r}'(u)\|(-du) = \int_a^b f(\mathbf{r}(u))\,\|\mathbf{r}'(u)\|\,du.
\]
This is exactly the original integral. The sign does not change because $ds = \|\mathbf{r}'(t)\|\,dt > 0$ always; the norm $\|\cdot\|$ absorbs any sign change from reversing direction.

%%% PROBLEM 17 %%%
\item \textbf{Solution.}
Parameterize: $x = r\cos t$, $y = r\sin t$, for $t\in[0,\pi/2]$.

$\mathbf{r}'(t) = (-r\sin t, r\cos t)$, so $\|\mathbf{r}'(t)\| = r$.

On $C$: $\sqrt{x^2+y^2} = \sqrt{r^2\cos^2 t + r^2\sin^2 t} = r$.
\[
\int_C \sqrt{x^2+y^2}\,ds = \int_0^{\pi/2} r\cdot r\,dt = r^2\int_0^{\pi/2}dt = \boxed{\frac{\pi r^2}{2}}.
\]

%%% PROBLEM 18 %%%
\item \textbf{Solution.}
\textbf{Segment 1} ($C_1$): from $(0,0)$ to $(3,0)$ along the $x$-axis. Parameterize: $\mathbf{r}_1(t) = (t,0)$, $0\le t\le 3$. Then $\|\mathbf{r}_1'(t)\|=1$, $y=0$.
\[
\int_{C_1}(2x+3y)\,ds = \int_0^3 2t\cdot 1\,dt = \left[t^2\right]_0^3 = 9.
\]

\textbf{Segment 2} ($C_2$): from $(3,0)$ to $(3,2)$ vertically. Parameterize: $\mathbf{r}_2(t) = (3,t)$, $0\le t\le 2$. Then $\|\mathbf{r}_2'(t)\|=1$, $x=3$.
\[
\int_{C_2}(2x+3y)\,ds = \int_0^2 (6+3t)\cdot 1\,dt = \left[6t+\tfrac{3}{2}t^2\right]_0^2 = 12+6 = 18.
\]

Total: $\displaystyle\int_C (2x+3y)\,ds = 9 + 18 = \boxed{27}$.

%%% PROBLEM 19 %%%
\item \textbf{Solution.}
Parameterize $C$: $\mathbf{r}(t) = (2t, 3t)$, $0\le t\le 1$. Then $\mathbf{r}'(t)=(2,3)$.

$\mathbf{F}(\mathbf{r}(t)) = (y,x) = (3t, 2t)$.
\[
\int_C \mathbf{F}\cdot d\mathbf{r} = \int_0^1 (3t,2t)\cdot(2,3)\,dt = \int_0^1(6t+6t)\,dt = \int_0^1 12t\,dt = 6t^2\Big|_0^1 = \boxed{6}.
\]

%%% PROBLEM 20 %%%
\item \textbf{Solution.}
$\mathbf{F}=(x-y, x^2)$, so $P=x-y$, $Q=x^2$. The rectangle has four edges traversed counterclockwise.

\textbf{$C_1$}: $(0,0)\to(2,0)$: $\mathbf{r}=(t,0)$, $t\in[0,2]$, $d\mathbf{r}=(1,0)dt$.
\[
\int_{C_1} = \int_0^2 (t-0)\cdot 1\,dt = \int_0^2 t\,dt = 2.
\]

\textbf{$C_2$}: $(2,0)\to(2,1)$: $\mathbf{r}=(2,t)$, $t\in[0,1]$, $d\mathbf{r}=(0,1)dt$.
\[
\int_{C_2} = \int_0^1 4\cdot 1\,dt = 4.
\]

\textbf{$C_3$}: $(2,1)\to(0,1)$: $\mathbf{r}=(2-t,1)$, $t\in[0,2]$, $d\mathbf{r}=(-1,0)dt$.
\[
\int_{C_3} = \int_0^2 ((2-t)-1)(-1)\,dt = \int_0^2 -(1-t)\,dt = \int_0^2 (t-1)\,dt = \left[\tfrac{t^2}{2}-t\right]_0^2 = 0.
\]

\textbf{$C_4$}: $(0,1)\to(0,0)$: $\mathbf{r}=(0,1-t)$, $t\in[0,1]$, $d\mathbf{r}=(0,-1)dt$.
\[
\int_{C_4} = \int_0^1 0\cdot(-1)\,dt = 0.
\]

Total: $2+4+0+0 = \boxed{6}$.

%%% PROBLEM 21 %%%
\item \textbf{Solution.}
$\mathbf{F}=(\ln(1+y^2), e^x)$, $C$: $x^2+y^2=4$ counterclockwise.

Parameterize: $\mathbf{r}(t) = (2\cos t, 2\sin t)$, $t\in[0,2\pi]$, $\mathbf{r}'(t)=(-2\sin t, 2\cos t)$.

By Green's Theorem:
\[
\oint_C \mathbf{F}\cdot d\mathbf{r} = \iint_D\left(\frac{\partial Q}{\partial x}-\frac{\partial P}{\partial y}\right)dA
\]
where $P = \ln(1+y^2)$, $Q = e^x$.
\[
\frac{\partial Q}{\partial x} = e^x, \qquad \frac{\partial P}{\partial y} = \frac{2y}{1+y^2}.
\]
\[
\oint_C \mathbf{F}\cdot d\mathbf{r} = \iint_D\left(e^x - \frac{2y}{1+y^2}\right)dA.
\]
The integrand $\frac{2y}{1+y^2}$ is an odd function in $y$ and $D$ is the disk $x^2+y^2\le 4$ which is symmetric about the $x$-axis, so:
\[
\iint_D \frac{2y}{1+y^2}\,dA = 0.
\]
Thus we need $\iint_D e^x\,dA$. Using polar coordinates $x=r\cos\theta$:
\[
\iint_D e^x\,dA = \int_0^{2\pi}\int_0^2 e^{r\cos\theta}\,r\,dr\,d\theta.
\]
This integral can be expressed in terms of modified Bessel functions. Using the identity $\frac{1}{2\pi}\int_0^{2\pi}e^{r\cos\theta}\,d\theta = I_0(r)$:
\[
= 2\pi\int_0^2 r\,I_0(r)\,dr.
\]
Since $\frac{d}{dr}[rI_1(r)] = rI_0(r)$, we get $2\pi[rI_1(r)]_0^2 = 4\pi I_1(2)$.

Therefore $\displaystyle\oint_C \mathbf{F}\cdot d\mathbf{r} = 4\pi I_1(2)$, where $I_1(2) \approx 1.5906$, giving approximately $\boxed{4\pi I_1(2) \approx 19.99}$.

Reversing to clockwise negates the integral: $\displaystyle\oint_{C^-}\mathbf{F}\cdot d\mathbf{r} = -4\pi I_1(2)$.

%%% PROBLEM 22 %%%
\item \textbf{Solution.}
If $C$ is piecewise smooth with segments $C_1$, $C_2$, $C_3$, then
\[
\int_C \mathbf{F}\cdot d\mathbf{r} = \int_{C_1}\mathbf{F}\cdot d\mathbf{r} + \int_{C_2}\mathbf{F}\cdot d\mathbf{r} + \int_{C_3}\mathbf{F}\cdot d\mathbf{r}.
\]
This follows from the additivity of integrals over adjoining curves: the endpoint of $C_1$ equals the starting point of $C_2$, and so on.

This strategy is often used in applications because each smooth segment may have a simple parameterization (e.g., a line segment or circular arc), even though the whole curve $C$ does not admit a single simple formula.

%%% PROBLEM 23 %%%
\item \textbf{Solution.}
Let $\mathbf{F} = (3y, -2x)$, so $P=3y$ and $Q=-2x$.

By Green's Theorem (for counterclockwise $C$ enclosing region $D$):
\[
\oint_C \mathbf{F}\cdot d\mathbf{r} = \iint_D\left(\frac{\partial Q}{\partial x} - \frac{\partial P}{\partial y}\right)dA = \iint_D(-2-3)\,dA = -5\,\text{Area}(D).
\]
This is nonzero for any region with positive area, confirming the integral depends on the loop.

If $C$ is traversed \textbf{clockwise}, the integral picks up a sign change:
\[
\oint_{C_{\text{cw}}} \mathbf{F}\cdot d\mathbf{r} = +5\,\text{Area}(D).
\]
So the clockwise result is the negative of the counterclockwise result.

%%% PROBLEM 24 %%%
\item \textbf{Solution.}
Along $x=1$, parameterize: $\mathbf{r}(t) = (1,t)$, $0\le t\le 2$, $\mathbf{r}'(t)=(0,1)$.

$\mathbf{F}(1,t) = (1+1, 2(1)(t)-1) = (2, 2t-1)$.
\[
\int_C \mathbf{F}\cdot d\mathbf{r} = \int_0^2 (2,2t-1)\cdot(0,1)\,dt = \int_0^2(2t-1)\,dt = \left[t^2-t\right]_0^2 = 4-2 = \boxed{2}.
\]

Reversing the path (from $(1,2)$ to $(1,0)$) gives $-2$. This is because for a vector line integral $\int_C \mathbf{F}\cdot d\mathbf{r}$, reversing the direction of traversal negates $d\mathbf{r}$, hence negates the integral.

%%% PROBLEM 25 %%%
\item \textbf{Solution.}
$\mathbf{F}(x,y) = (y,x)$, so $P=y$, $Q=x$.
\[
\frac{\partial P}{\partial y} = 1, \qquad \frac{\partial Q}{\partial x} = 1.
\]
Since $\dfrac{\partial P}{\partial y} = \dfrac{\partial Q}{\partial x}$ everywhere on $\mathbb{R}^2$ (which is simply connected), $\mathbf{F}$ \textbf{is conservative}.

A potential function: $f_x = y \Rightarrow f = xy + g(y)$. Then $f_y = x + g'(y) = x$, so $g'(y)=0$, giving $g(y)=C$. Thus $f(x,y)=xy$.

%%% PROBLEM 26 %%%
\item \textbf{Solution.}
$\mathbf{F} = \left(\dfrac{-y}{x^2+y^2},\;\dfrac{x}{x^2+y^2}\right)$ on $\mathbb{R}^2\setminus\{(0,0)\}$.

Parameterize $C$: $x = r\cos t$, $y = r\sin t$, $0\le t\le 2\pi$.

$dx = -r\sin t\,dt$, $dy = r\cos t\,dt$.

On $C$: $x^2+y^2 = r^2$.
\begin{align*}
\oint_C \mathbf{F}\cdot d\mathbf{r} &= \int_0^{2\pi}\left[\frac{-r\sin t}{r^2}(-r\sin t) + \frac{r\cos t}{r^2}(r\cos t)\right]dt \\
&= \int_0^{2\pi}\left[\frac{\sin^2 t}{1} + \frac{\cos^2 t}{1}\right]dt = \int_0^{2\pi}1\,dt = \boxed{2\pi}.
\end{align*}

Although one can check $\frac{\partial P}{\partial y} = \frac{\partial Q}{\partial x}$ (both equal $\frac{y^2-x^2}{(x^2+y^2)^2}$), the field is \textbf{not conservative} because the domain $\mathbb{R}^2\setminus\{(0,0)\}$ is \emph{not simply connected}. The circulation around the origin is $2\pi \neq 0$, which is impossible for a conservative field.

%%% PROBLEM 27 %%%
\item \textbf{Solution.}
$\mathbf{G}=(x+\sin y,\;\cos x + y)$. Check: $P = x+\sin y$, $Q = \cos x + y$.
\[
\frac{\partial P}{\partial y} = \cos y, \qquad \frac{\partial Q}{\partial x} = -\sin x.
\]
Since $\cos y \neq -\sin x$ in general (e.g., at $(0,0)$: $1\neq 0$), $\mathbf{G}$ is \textbf{not conservative}.

Therefore we cannot use the potential function method. The integral $\int_C \mathbf{G}\cdot d\mathbf{r}$ from $(0,0)$ to $(2,1)$ is path-dependent and requires specifying a particular path.

%%% PROBLEM 28 %%%
\item \textbf{Solution.}
$\mathbf{F} = (yz, zx, xy)$.

\textbf{$C_1$}: $(0,0,0)\to(1,0,0)$. $\mathbf{r}(t)=(t,0,0)$, $t\in[0,1]$, $\mathbf{r}'=(1,0,0)$.
$\mathbf{F}=(0,0,0)$. So $\int_{C_1}=0$.

\textbf{$C_2$}: $(1,0,0)\to(1,2,0)$. $\mathbf{r}(t)=(1,t,0)$, $t\in[0,2]$, $\mathbf{r}'=(0,1,0)$.
$\mathbf{F}=(0,0,t)$. So $\mathbf{F}\cdot\mathbf{r}' = 0$, and $\int_{C_2}=0$.

\textbf{$C_3$}: $(1,2,0)\to(1,2,3)$. $\mathbf{r}(t)=(1,2,t)$, $t\in[0,3]$, $\mathbf{r}'=(0,0,1)$.
$\mathbf{F}=(2t, t, 2)$. So $\mathbf{F}\cdot\mathbf{r}' = 2$, and $\int_{C_3} = \int_0^3 2\,dt = 6$.

Total: $0+0+6 = \boxed{6}$.

To check if $\mathbf{F}$ is conservative, check $\nabla\times\mathbf{F}$:
\[
\nabla\times\mathbf{F} = \left(\frac{\partial(xy)}{\partial y}-\frac{\partial(zx)}{\partial z},\;\frac{\partial(yz)}{\partial z}-\frac{\partial(xy)}{\partial x},\;\frac{\partial(zx)}{\partial x}-\frac{\partial(yz)}{\partial y}\right) = (x-x, y-y, z-z) = \mathbf{0}.
\]
Since the curl is zero on all of $\mathbb{R}^3$ (simply connected), $\mathbf{F}$ \textbf{is conservative}. A potential function is $f(x,y,z)=xyz$. We verify: $f(1,2,3)-f(0,0,0)=6-0=6$.

%%% PROBLEM 29 %%%
\item \textbf{Solution.}
\textbf{Green's Theorem} states: if $C$ is a positively oriented (counterclockwise), piecewise smooth, simple closed curve bounding a region $D$, and $P,Q$ have continuous partial derivatives on an open set containing $D$, then
\[
\oint_C P\,dx + Q\,dy = \iint_D \left(\frac{\partial Q}{\partial x} - \frac{\partial P}{\partial y}\right)dA.
\]

\textbf{Conceptual explanation:} The line integral measures the total ``circulation'' of the vector field $(P,Q)$ around the boundary. Green's Theorem says this equals the sum of all the ``local rotation'' ($\partial Q/\partial x - \partial P/\partial y$) accumulated over the interior region. Intuitively, if we partition $D$ into infinitesimally small rectangles and sum the circulation around each, all interior edge contributions cancel (shared edges are traversed in opposite directions by adjacent rectangles), leaving only the outer boundary $C$. The double integral simply adds up the microscopic rotation at every interior point.

%%% PROBLEM 30 %%%
\item \textbf{Solution.}
$P = xy$, $Q = x^2$, rectangle $[0,2]\times[0,1]$.

\textbf{Direct computation:}

$C_1$: $(0,0)\to(2,0)$, $y=0$: $\int_0^2 (0)(1)\,dt + 0 = 0$.

$C_2$: $(2,0)\to(2,1)$, $x=2$: $\mathbf{r}=(2,t)$, $dx=0$, $dy=dt$.
$\int_0^1 4\,dt = 4$.

$C_3$: $(2,1)\to(0,1)$, $y=1$: $\mathbf{r}=(2-t,1)$, $dx=-dt$, $dy=0$.
$\int_0^2 (2-t)(1)(-1)\,dt = -\int_0^2(2-t)\,dt = -(2\cdot2-2)=-2$.

$C_4$: $(0,1)\to(0,0)$, $x=0$: $\int_0^1 0\,dt = 0$.

Direct total: $0+4-2+0 = 2$.

\textbf{Green's Theorem:}
\[
\frac{\partial Q}{\partial x} - \frac{\partial P}{\partial y} = 2x - x = x.
\]
\[
\iint_D x\,dA = \int_0^1\int_0^2 x\,dx\,dy = \int_0^1 2\,dy = 2.
\]

$\boxed{2}$.

%%% PROBLEM 31 %%%
\item \textbf{Solution.}
$P=-y$, $Q=x$, $C$: unit circle counterclockwise, $D$: unit disk.
\[
\frac{\partial Q}{\partial x}-\frac{\partial P}{\partial y} = 1-(-1)=2.
\]
By Green's Theorem:
\[
\oint_C \mathbf{F}\cdot d\mathbf{r} = \iint_D 2\,dA = 2\cdot\text{Area}(D) = 2\pi(1)^2 = \boxed{2\pi}.
\]

\textbf{Physical interpretation:} This integral measures the \emph{circulation} of $\mathbf{F}$ around the unit circle. The field $\mathbf{F}=(-y,x)$ rotates counterclockwise with speed equal to the distance from the origin. The constant ``curl'' of $2$ means uniform rotation everywhere, so the total circulation around the unit circle is $2\pi$---consistent with a rigid-body rotation at angular velocity $1$.

%%% PROBLEM 32 %%%
\item \textbf{Solution.}
The region $D$ is bounded by $y = x^2$ (below) and $y=4$ (above), so $x$ ranges from $-2$ to $2$.

By Green's Theorem (in reverse):
\[
\iint_D\left(\frac{\partial Q}{\partial x}-\frac{\partial P}{\partial y}\right)dA = \oint_C P\,dx + Q\,dy,
\]
where $C = \partial D$ is traversed counterclockwise.

The boundary consists of two pieces:
\begin{itemize}
\item $C_1$: along $y = x^2$ from $(-2,4)$ to $(2,4)$ (bottom, left to right).
\item $C_2$: along $y = 4$ from $(2,4)$ to $(-2,4)$ (top, right to left).
\end{itemize}

The resulting line integral is:
\[
\oint_C P\,dx + Q\,dy = \int_{C_1} P\,dx + Q\,dy + \int_{C_2} P\,dx + Q\,dy.
\]

One advantage of switching to the line integral form is that if $P$ and $Q$ are simple functions, a one-dimensional boundary integral can be easier to evaluate than a two-dimensional area integral, especially when the boundary has a convenient parameterization.

%%% PROBLEM 33 %%%
\item \textbf{Solution.}
When we partition $D$ into small rectangles and sum the line integrals around each sub-rectangle, every \emph{interior} edge is shared by exactly two adjacent sub-rectangles. These two rectangles traverse the shared edge in \emph{opposite directions}: what is the right edge of one rectangle is the left edge of its neighbor, and the orientations are opposite. Thus the contributions from interior edges cancel in pairs.

Only the edges lying on the \emph{outer boundary} of $D$ belong to a single sub-rectangle and have no partner to cancel with. These surviving contributions sum to exactly the line integral around the outer boundary $\partial D$. This is the discrete version of Green's Theorem: the sum of local circulations (double integral of curl) equals the global circulation (boundary line integral).

%%% PROBLEM 34 %%%
\item \textbf{Solution.}
$P = x^2y^2$, $Q = xy^3$. Triangle with vertices $(0,0)$, $(2,0)$, $(0,2)$.
\[
\frac{\partial Q}{\partial x} = y^3, \qquad \frac{\partial P}{\partial y} = 2x^2y.
\]
\[
\frac{\partial Q}{\partial x}-\frac{\partial P}{\partial y} = y^3 - 2x^2y.
\]
By Green's Theorem:
\[
\oint_C = \iint_D (y^3 - 2x^2y)\,dA.
\]
The triangle: $0\le x \le 2$, $0\le y\le 2-x$.
\begin{align*}
&= \int_0^2\int_0^{2-x}(y^3-2x^2y)\,dy\,dx \\
&= \int_0^2\left[\frac{y^4}{4}-x^2y^2\right]_0^{2-x}dx \\
&= \int_0^2\left[\frac{(2-x)^4}{4}-x^2(2-x)^2\right]dx.
\end{align*}
Let $u = 2-x$:
\[
\frac{(2-x)^4}{4} - x^2(2-x)^2 = \frac{u^4}{4}-(2-u)^2 u^2.
\]
Expand $(2-u)^2 u^2 = (4-4u+u^2)u^2 = 4u^2 - 4u^3 + u^4$.
\[
= \frac{u^4}{4}-4u^2+4u^3-u^4 = -\frac{3u^4}{4}+4u^3-4u^2.
\]
Back in $x$ with $u=2-x$, $du=-dx$:
\[
\int_0^2\left[-\frac{3(2-x)^4}{4}+4(2-x)^3-4(2-x)^2\right]dx.
\]
Let $u=2-x$:
\[
= \int_0^2\left[-\frac{3u^4}{4}+4u^3-4u^2\right]du = \left[-\frac{3u^5}{20}+u^4-\frac{4u^3}{3}\right]_0^2
\]
\[
= -\frac{3(32)}{20}+16-\frac{32}{3} = -\frac{96}{20}+16-\frac{32}{3} = -\frac{24}{5}+16-\frac{32}{3}.
\]
Common denominator $15$:
\[
= \frac{-72+240-160}{15} = \frac{8}{15}.
\]

$\boxed{\dfrac{8}{15}}$.

%%% PROBLEM 35 %%%
\item \textbf{Solution.}
$\mathbf{F}=(y,-x)$, so $P=y$, $Q=-x$.
\[
\frac{\partial Q}{\partial x}-\frac{\partial P}{\partial y} = -1 - 1 = -2.
\]
By Green's Theorem:
\[
\oint_C \mathbf{F}\cdot d\mathbf{r} = \iint_D(-2)\,dA = -2\,\text{Area}(D).
\]
Therefore the line integral is proportional to the area of the enclosed region, with proportionality constant $-2$:
\[
\oint_C (y\,dx - x\,dy) = -2\,\text{Area}(D). \qquad\blacksquare
\]

%%% PROBLEM 36 %%%
\item \textbf{Solution.}
$\mathbf{F}=(e^y, \sin x)$, so $P=e^y$, $Q=\sin x$.
\[
\frac{\partial Q}{\partial x} = \cos x, \qquad \frac{\partial P}{\partial y} = e^y.
\]
By Green's Theorem:
\[
\oint_C \mathbf{F}\cdot d\mathbf{r} = \iint_D(\cos x - e^y)\,dA,
\]
where $D$ is the interior of the ellipse $\frac{x^2}{4}+\frac{y^2}{9}=1$.

Use the substitution $x = 2u$, $y = 3v$ (mapping the ellipse to the unit disk $u^2+v^2\le 1$), with Jacobian $|J|=6$:
\[
= \iint_{u^2+v^2\le 1}(\cos(2u)-e^{3v})\cdot 6\,du\,dv.
\]
By symmetry, $\cos(2u)$ is even in both $u$ and $v$, and $e^{3v}$ involves only $v$.

For $\iint \cos(2u)\,du\,dv$ over the unit disk, switch to polar: $u=r\cos\theta$:
\[
\int_0^{2\pi}\int_0^1\cos(2r\cos\theta)\,r\,dr\,d\theta = 2\pi\int_0^1 r\,J_0(2r)\,dr,
\]
where $J_0$ is the Bessel function. Using $\frac{d}{dr}[rJ_1(2r)\cdot\frac{1}{2}]=\frac{1}{2}\cdot[J_1(2r)+2rJ_0(2r)\cdot\frac{1}{2}]$... this becomes intricate.

A more direct approach: use the identity $\int_0^1 r J_0(2r)\,dr = \frac{J_1(2)}{2}$, so $\iint\cos(2u)\,du\,dv = 2\pi\cdot\frac{J_1(2)}{2} = \pi J_1(2)$.

For $\iint e^{3v}\,du\,dv$ over the unit disk: $v = r\sin\theta$:
\[
\int_0^{2\pi}\int_0^1 e^{3r\sin\theta}\,r\,dr\,d\theta = 2\pi\int_0^1 r\,I_0(3r)\,dr = 2\pi\cdot\frac{I_1(3)}{3},
\]
using $\int_0^1 r I_0(3r)\,dr = \frac{I_1(3)}{3}$.

Therefore:
\[
\oint_C \mathbf{F}\cdot d\mathbf{r} = 6\left[\pi J_1(2) - \frac{2\pi I_1(3)}{3}\right] = 6\pi J_1(2) - 4\pi I_1(3).
\]

Using $J_1(2)\approx 0.5767$ and $I_1(3)\approx 3.9534$:
\[
\approx 6\pi(0.5767)-4\pi(3.9534) \approx 10.87 - 49.68 \approx \boxed{-38.81}.
\]

%%% PROBLEM 37 %%%
\item \textbf{Solution.}
Green's Theorem states (for counterclockwise $C$):
\[
\oint_{C_{\text{ccw}}} P\,dx+Q\,dy = \iint_D\left(\frac{\partial Q}{\partial x}-\frac{\partial P}{\partial y}\right)dA.
\]
When we reverse the orientation from counterclockwise to clockwise, $d\mathbf{r}$ is negated (the tangent vector reverses), so:
\[
\oint_{C_{\text{cw}}} P\,dx+Q\,dy = -\oint_{C_{\text{ccw}}} P\,dx+Q\,dy = -\iint_D\left(\frac{\partial Q}{\partial x}-\frac{\partial P}{\partial y}\right)dA.
\]
The double integral $\iint_D(\partial Q/\partial x - \partial P/\partial y)\,dA$ does not depend on orientation---it is a property of the region and the field. The sign convention (positive = counterclockwise) is what links the line integral to the double integral.

%%% PROBLEM 38 %%%
\item \textbf{Solution.}
$\mathbf{F} = \left(\frac{-y}{x^2+y^2},\;\frac{x}{x^2+y^2}\right)$ on $\mathbb{R}^2\setminus\{(0,0)\}$.

First check:
\[
\frac{\partial Q}{\partial x} = \frac{(x^2+y^2)-x(2x)}{(x^2+y^2)^2} = \frac{y^2-x^2}{(x^2+y^2)^2},
\]
\[
\frac{\partial P}{\partial y} = \frac{-(x^2+y^2)+y(2y)}{(x^2+y^2)^2} = \frac{y^2-x^2}{(x^2+y^2)^2}.
\]
So $\frac{\partial Q}{\partial x}-\frac{\partial P}{\partial y}=0$ wherever $\mathbf{F}$ is defined.

If Green's Theorem applied directly, we would conclude $\oint_C \mathbf{F}\cdot d\mathbf{r}=\iint_D 0\,dA = 0$. However, $\mathbf{F}$ is undefined at the origin, which lies inside $C$ (the circle of radius $r$). The domain is \textbf{not simply connected}, so Green's Theorem does not apply to the region enclosed by $C$.

Computing directly with parameterization $x=r\cos t$, $y=r\sin t$, $t\in[0,2\pi]$:
\begin{align*}
\oint_C \mathbf{F}\cdot d\mathbf{r} &= \int_0^{2\pi}\left[\frac{-r\sin t}{r^2}(-r\sin t)+\frac{r\cos t}{r^2}(r\cos t)\right]dt \\
&= \int_0^{2\pi}(\sin^2 t + \cos^2 t)\,dt = \boxed{2\pi}.
\end{align*}

This demonstrates the crucial role of simple connectivity: $\partial Q/\partial x = \partial P/\partial y$ everywhere in the domain, yet the field is not conservative because the domain has a ``hole'' at the origin.

%%% PROBLEM 39 %%%
\item \textbf{Solution.}
$\mathbf{F}=(x+2y,\;x-y)$, so $P=x+2y$, $Q=x-y$.

\[
\text{div}(\mathbf{F}) = \frac{\partial P}{\partial x}+\frac{\partial Q}{\partial y} = 1+(-1) = \boxed{0}.
\]
\[
\text{curl}(\mathbf{F})\text{ (2D)} = \frac{\partial Q}{\partial x}-\frac{\partial P}{\partial y} = 1-2 = \boxed{-1}.
\]

The \textbf{divergence} indicates expansion/compression (here zero, meaning the field is area-preserving). The \textbf{curl} indicates rotation (here $-1$, meaning clockwise rotation).

%%% PROBLEM 40 %%%
\item \textbf{Solution.}
$\mathbf{F}=(x^2, 2y, -3z)$.

\[
\text{div}(\mathbf{F}) = \frac{\partial}{\partial x}(x^2)+\frac{\partial}{\partial y}(2y)+\frac{\partial}{\partial z}(-3z) = 2x+2-3 = 2x-1.
\]

This is not identically zero, so $\mathbf{F}$ is \textbf{not incompressible} (not divergence-free).

\[
\text{curl}(\mathbf{F}) = \begin{vmatrix}\mathbf{i}&\mathbf{j}&\mathbf{k}\\ \partial_x&\partial_y&\partial_z\\ x^2&2y&-3z\end{vmatrix} = \mathbf{i}(0-0)-\mathbf{j}(0-0)+\mathbf{k}(0-0) = \boxed{\mathbf{0}}.
\]
The curl is zero, so $\mathbf{F}$ is irrotational.

%%% PROBLEM 41 %%%
\item \textbf{Solution.}
If $\mathbf{F}=\nabla f = \left(\frac{\partial f}{\partial x},\frac{\partial f}{\partial y},\frac{\partial f}{\partial z}\right)$, then
\[
\nabla\times\mathbf{F} = \begin{vmatrix}\mathbf{i}&\mathbf{j}&\mathbf{k}\\ \partial_x&\partial_y&\partial_z\\ f_x & f_y & f_z\end{vmatrix}.
\]
The $\mathbf{i}$-component is $\frac{\partial f_z}{\partial y}-\frac{\partial f_y}{\partial z} = f_{zy}-f_{yz}$. By Clairaut's theorem (assuming $f$ has continuous second partial derivatives), $f_{zy}=f_{yz}$, so this component is $0$.

Similarly, the $\mathbf{j}$-component is $f_{xz}-f_{zx}=0$ and the $\mathbf{k}$-component is $f_{yx}-f_{xy}=0$.

Therefore $\nabla\times(\nabla f) = \mathbf{0}$ for any sufficiently smooth scalar function $f$. Every conservative field is irrotational. $\blacksquare$

%%% PROBLEM 42 %%%
\item \textbf{Solution.}
$\mathbf{F}=(e^x\sin y,\;e^y\sin z,\;e^z\sin x)$.
\[
\text{div}(\mathbf{F}) = \frac{\partial}{\partial x}(e^x\sin y)+\frac{\partial}{\partial y}(e^y\sin z)+\frac{\partial}{\partial z}(e^z\sin x)
\]
\[
= e^x\sin y + e^y\sin z + e^z\sin x.
\]
This is \textbf{not identically zero}. For example, at $(\pi/2, \pi/2, \pi/2)$:
\[
\text{div}(\mathbf{F}) = e^{\pi/2}\sin(\pi/2)+e^{\pi/2}\sin(\pi/2)+e^{\pi/2}\sin(\pi/2) = 3e^{\pi/2}\neq 0.
\]
Therefore $\mathbf{F}$ is \textbf{not divergence-free}.

%%% PROBLEM 43 %%%
\item \textbf{Solution.}
$\mathbf{F}=(x^3, y^3)$.
\[
\text{div}(\mathbf{F}) = 3x^2 + 3y^2 = 3(x^2+y^2).
\]
This equals zero only at the origin $(0,0)$ and is strictly positive everywhere else.

\textbf{Physical interpretation:} Since $\text{div}(\mathbf{F})=3(x^2+y^2)>0$ for $(x,y)\neq(0,0)$, the field acts as a \emph{source} at every non-origin point: fluid is being ``created'' or expanding outward. The source strength increases with distance from the origin. Only at the origin itself does the divergence vanish (neither source nor sink there).

%%% PROBLEM 44 %%%
\item \textbf{Solution.}
\begin{enumerate}
\item $\mathbf{F}_1=(x,-y)$:
\[
\text{div}(\mathbf{F}_1)=1+(-1)=0, \quad \text{curl}(\mathbf{F}_1)=\frac{\partial(-y)}{\partial x}-\frac{\partial(x)}{\partial y}=0-0=0.
\]
At the origin, both divergence and curl are zero. This is a \textbf{saddle} (hyperbolic flow), not primarily source, sink, or rotation.

\item $\mathbf{F}_2=(2x,2y)$:
\[
\text{div}(\mathbf{F}_2)=2+2=4>0, \quad \text{curl}(\mathbf{F}_2)=0-0=0.
\]
Divergence is positive, curl is zero. The origin is a \textbf{source}.

\item $\mathbf{F}_3=(-y,x)$:
\[
\text{div}(\mathbf{F}_3)=0+0=0, \quad \text{curl}(\mathbf{F}_3)=1-(-1)=2.
\]
Divergence is zero, curl is nonzero. The origin is a \textbf{rotation} (counterclockwise).
\end{enumerate}

%%% PROBLEM 45 %%%
\item \textbf{Solution.}
The identity $\nabla\cdot(\nabla\times\mathbf{F})=0$ holds for any smooth vector field $\mathbf{F}$.

If the velocity field $\mathbf{v}$ of a fluid is purely rotational, meaning $\mathbf{v}=\nabla\times\mathbf{A}$ for some vector potential $\mathbf{A}$, then
\[
\nabla\cdot\mathbf{v} = \nabla\cdot(\nabla\times\mathbf{A})=0.
\]
This means the divergence of $\mathbf{v}$ is zero everywhere.

By the Divergence Theorem, for any closed surface $S$ bounding a volume $V$:
\[
\oiint_S \mathbf{v}\cdot d\mathbf{S} = \iiint_V\nabla\cdot\mathbf{v}\,dV = 0.
\]
The net outward flux through any closed surface is zero, meaning no fluid is created or destroyed inside $V$. In a purely swirling (rotational) motion, the fluid merely circulates without any net outflow or inflow---mass is conserved. $\blacksquare$

%%% PROBLEM 46 %%%
\item \textbf{Solution.}
$\mathbf{F}=(yz, zx, xy)$.
\[
\text{div}(\mathbf{F}) = \frac{\partial(yz)}{\partial x}+\frac{\partial(zx)}{\partial y}+\frac{\partial(xy)}{\partial z} = 0+0+0 = \boxed{0}.
\]

\[
\text{curl}(\mathbf{F}) = \begin{vmatrix}\mathbf{i}&\mathbf{j}&\mathbf{k}\\ \partial_x&\partial_y&\partial_z\\ yz&zx&xy\end{vmatrix}
\]
\[
= \mathbf{i}(x-x)-\mathbf{j}(y-y)+\mathbf{k}(z-z) = \boxed{\mathbf{0}}.
\]

Both divergence and curl are zero. The divergence being zero means no net creation or destruction of fluid (incompressible). The curl being zero means no rotation (irrotational). In fact, $\mathbf{F}=\nabla(xyz)$, confirming it is conservative.

\emph{Note:} The problem statement suggests the curl is nonzero, but computation shows both are zero. The field is both divergence-free and irrotational.

%%% PROBLEM 47 %%%
\item \textbf{Solution.}
In cylindrical coordinates with $\mathbf{v}=(v_r, v_\theta, v_z)=(0,\alpha r, 0)$:

\textbf{Divergence} in cylindrical coordinates:
\[
\nabla\cdot\mathbf{v} = \frac{1}{r}\frac{\partial(rv_r)}{\partial r}+\frac{1}{r}\frac{\partial v_\theta}{\partial\theta}+\frac{\partial v_z}{\partial z} = \frac{1}{r}\frac{\partial(0)}{\partial r}+\frac{1}{r}\frac{\partial(\alpha r)}{\partial\theta}+0 = 0+0+0 = \boxed{0}.
\]
The flow is \textbf{incompressible}.

\textbf{Curl} in cylindrical coordinates:
\[
(\nabla\times\mathbf{v})_r = \frac{1}{r}\frac{\partial v_z}{\partial\theta}-\frac{\partial v_\theta}{\partial z} = 0-0=0,
\]
\[
(\nabla\times\mathbf{v})_\theta = \frac{\partial v_r}{\partial z}-\frac{\partial v_z}{\partial r} = 0-0=0,
\]
\[
(\nabla\times\mathbf{v})_z = \frac{1}{r}\frac{\partial(rv_\theta)}{\partial r}-\frac{1}{r}\frac{\partial v_r}{\partial\theta} = \frac{1}{r}\frac{\partial(\alpha r^2)}{\partial r}-0 = \frac{2\alpha r}{r} = 2\alpha.
\]
\[
\nabla\times\mathbf{v} = \boxed{2\alpha\,\mathbf{e}_z}.
\]

The curl is nonzero (pointing along $z$), confirming the flow \textbf{rotates about the $z$-axis}. This is solid-body rotation with angular velocity $\alpha$.

%%% PROBLEM 48 %%%
\item \textbf{Solution.}
In spherical coordinates with $\mathbf{F}=(f(r),0,0)$ (purely radial), the divergence is:
\[
\nabla\cdot\mathbf{F} = \frac{1}{r^2}\frac{d}{dr}\bigl(r^2 f(r)\bigr).
\]

Expanding:
\[
\nabla\cdot\mathbf{F} = \frac{1}{r^2}\left[2r\,f(r)+r^2 f'(r)\right] = \frac{2f(r)}{r}+f'(r).
\]

For this to be zero for $r>0$:
\[
\frac{d}{dr}\bigl(r^2 f(r)\bigr)=0 \implies r^2 f(r)=C \implies \boxed{f(r)=\frac{C}{r^2}}
\]
for some constant $C$. This is the inverse-square law (e.g., gravitational or electrostatic fields).

%%% PROBLEM 49 %%%
\item \textbf{Solution.}
Let $\mathbf{F}=(F_1,F_2,F_3)$ and $f$ be a scalar function.

\textbf{Left side:}
\[
\nabla\cdot(f\mathbf{F}) = \frac{\partial(fF_1)}{\partial x}+\frac{\partial(fF_2)}{\partial y}+\frac{\partial(fF_3)}{\partial z}.
\]
By the product rule for each term:
\[
= \left(f\frac{\partial F_1}{\partial x}+F_1\frac{\partial f}{\partial x}\right)+\left(f\frac{\partial F_2}{\partial y}+F_2\frac{\partial f}{\partial y}\right)+\left(f\frac{\partial F_3}{\partial z}+F_3\frac{\partial f}{\partial z}\right).
\]

Regrouping:
\[
= f\left(\frac{\partial F_1}{\partial x}+\frac{\partial F_2}{\partial y}+\frac{\partial F_3}{\partial z}\right)+\left(F_1\frac{\partial f}{\partial x}+F_2\frac{\partial f}{\partial y}+F_3\frac{\partial f}{\partial z}\right).
\]

\textbf{Right side:}
\[
f(\nabla\cdot\mathbf{F})+\mathbf{F}\cdot(\nabla f) = f\left(\frac{\partial F_1}{\partial x}+\frac{\partial F_2}{\partial y}+\frac{\partial F_3}{\partial z}\right)+\left(F_1 f_x+F_2 f_y+F_3 f_z\right).
\]

Both sides are identical. $\blacksquare$

\textbf{Physical interpretation:} The net outflow of the quantity $f\mathbf{F}$ has two contributions: (1) $f\,\nabla\cdot\mathbf{F}$, which is the outflow due to $\mathbf{F}$ itself expanding/contracting, weighted by $f$; and (2) $\mathbf{F}\cdot\nabla f$, which is the transport of $f$ due to the flow $\mathbf{F}$ carrying $f$ from regions of high $f$ to low $f$ (or vice versa).

%%% PROBLEM 50 %%%
\item \textbf{Solution.}
\[
\nabla\times(\nabla f) = \nabla\times\left(\frac{\partial f}{\partial x},\frac{\partial f}{\partial y},\frac{\partial f}{\partial z}\right).
\]
\[
= \left(\frac{\partial^2 f}{\partial y\partial z}-\frac{\partial^2 f}{\partial z\partial y}\right)\mathbf{i} - \left(\frac{\partial^2 f}{\partial x\partial z}-\frac{\partial^2 f}{\partial z\partial x}\right)\mathbf{j} + \left(\frac{\partial^2 f}{\partial x\partial y}-\frac{\partial^2 f}{\partial y\partial x}\right)\mathbf{k}.
\]
By Clairaut's theorem ($f_{yz}=f_{zy}$, etc.), every component is zero:
\[
\nabla\times(\nabla f)=\mathbf{0}. \qquad\blacksquare
\]

This means any gradient field is \textbf{irrotational}: there is no circulation or rotation.

\textbf{Physical application:} The electrostatic field $\mathbf{E}=-\nabla V$ (where $V$ is the electric potential) is irrotational: $\nabla\times\mathbf{E}=\mathbf{0}$ in the static case. This is one of Maxwell's equations for electrostatics.

%%% PROBLEM 51 %%%
\item \textbf{Solution.}
If $\nabla\cdot\mathbf{F}=3$ everywhere, then at every point, the field has a positive outflow rate of $3$ per unit volume. This means fluid is being ``created'' uniformly throughout space at a constant rate---$\mathbf{F}$ behaves like a \textbf{uniform source of fluid}.

To verify, use the Divergence Theorem on a sphere $S$ of radius $R$:
\[
\oiint_S \mathbf{F}\cdot\mathbf{n}\,dS = \iiint_V \nabla\cdot\mathbf{F}\,dV = \iiint_V 3\,dV = 3\cdot\frac{4}{3}\pi R^3 = 4\pi R^3.
\]
The total outward flux through the sphere is $4\pi R^3$, which grows with the volume, confirming the uniform source interpretation.

%%% PROBLEM 52 %%%
\item \textbf{Solution.}
$\mathbf{F}=(x^2, y^2, z^2)$.

\textbf{Step 1:} Compute $\nabla\times\mathbf{F}$:
\[
\nabla\times\mathbf{F} = (0-0)\mathbf{i}-(0-0)\mathbf{j}+(0-0)\mathbf{k} = \mathbf{0}.
\]

\textbf{Step 2:} $\nabla\times(\nabla\times\mathbf{F}) = \nabla\times\mathbf{0} = \mathbf{0}$.

\textbf{Step 3:} Verify with the identity $\nabla(\nabla\cdot\mathbf{F})-\nabla^2\mathbf{F}$:

$\nabla\cdot\mathbf{F} = 2x+2y+2z$.

$\nabla(\nabla\cdot\mathbf{F}) = \nabla(2x+2y+2z)=(2,2,2)$.

$\nabla^2\mathbf{F} = (\nabla^2(x^2),\;\nabla^2(y^2),\;\nabla^2(z^2)) = (2,2,2)$.

Therefore:
\[
\nabla(\nabla\cdot\mathbf{F})-\nabla^2\mathbf{F} = (2,2,2)-(2,2,2)=\mathbf{0}.
\]
Both sides agree: $\nabla\times(\nabla\times\mathbf{F})=\mathbf{0}$.

%%% PROBLEM 53 %%%
\item \textbf{Solution.}
\begin{enumerate}
\item $\mathbf{F}=\nabla(x^2+2y^2) = (2x, 4y)$, so $P=2x$, $Q=4y$.

\item $\text{curl}(\mathbf{F}) = \frac{\partial Q}{\partial x}-\frac{\partial P}{\partial y} = 0 - 0 = \boxed{0}$.

$\text{div}(\mathbf{F}) = \frac{\partial P}{\partial x}+\frac{\partial Q}{\partial y} = 2+4 = \boxed{6}$.

\item Yes, $\text{curl}(\mathbf{F})=0$ confirms that gradient fields are always irrotational, consistent with the identity $\nabla\times(\nabla f)=\mathbf{0}$.
\end{enumerate}

%%% PROBLEM 54 %%%
\item \textbf{Solution.}
The surface $x^2+y^2=4z^2$ is a cone. Rewrite as $z = \pm\frac{1}{2}\sqrt{x^2+y^2}$.

Using cylindrical-style parameters, let $x=r\cos\theta$, $y=r\sin\theta$, and then $z=\pm\frac{r}{2}$, where $r\ge 0$.

\textbf{Parametric form:}
\[
\mathbf{r}(r,\theta) = \left(r\cos\theta,\; r\sin\theta,\; \pm\frac{r}{2}\right),
\]
with parameter constraints: $r\ge 0$, $0\le\theta\le 2\pi$ (or $0\le\theta <2\pi$).

The $+$ sign gives the upper cone ($z\ge 0$) and the $-$ sign gives the lower cone ($z\le 0$). Together they form the full double cone.

%%% PROBLEM 55 %%%
\item \textbf{Solution.}
The surface is $x = y^2+2z^2$ with $x\le 4$.

Let $y = u$ and $z = v$ be the parameters. Then $x = u^2+2v^2$.

\textbf{Parametrization:}
\[
\mathbf{r}(u,v)=(u^2+2v^2,\; u,\; v).
\]

The constraint $x\le 4$ becomes $u^2+2v^2\le 4$, which is an elliptical region in the $(u,v)$-plane:
\[
\frac{u^2}{4}+\frac{v^2}{2}\le 1.
\]
So the domain is: $\{(u,v): u^2+2v^2\le 4\}$, i.e., $-2\le u\le 2$ and $-\sqrt{(4-u^2)/2}\le v\le\sqrt{(4-u^2)/2}$.

%%% PROBLEM 56 %%%
\item \textbf{Solution.}
$\mathbf{r}(u,v)=(u, v, \sqrt{u^2+v^2})$.

\begin{enumerate}
\item
\[
\mathbf{r}_u = \left(1, 0, \frac{u}{\sqrt{u^2+v^2}}\right), \qquad \mathbf{r}_v = \left(0, 1, \frac{v}{\sqrt{u^2+v^2}}\right).
\]

\item Let $s=\sqrt{u^2+v^2}$.
\[
\mathbf{n}=\mathbf{r}_u\times\mathbf{r}_v = \begin{vmatrix}\mathbf{i}&\mathbf{j}&\mathbf{k}\\ 1&0&u/s\\ 0&1&v/s\end{vmatrix} = \left(-\frac{u}{s},\;-\frac{v}{s},\;1\right).
\]

\item
\[
\|\mathbf{n}\| = \sqrt{\frac{u^2}{s^2}+\frac{v^2}{s^2}+1} = \sqrt{\frac{u^2+v^2}{u^2+v^2}+1} = \sqrt{1+1}=\boxed{\sqrt{2}}.
\]
\end{enumerate}

%%% PROBLEM 57 %%%
\item \textbf{Solution.}
The paraboloid is $z=9-x^2-y^2$ and the region is $x^2+y^2\le 5$.

For a surface $z=f(x,y)$, the surface area element is $dS = \sqrt{1+f_x^2+f_y^2}\,dA$.

Here $f_x=-2x$, $f_y=-2y$, so $\sqrt{1+4x^2+4y^2}$.

\[
\text{Surface Area} = \iint_{x^2+y^2\le 5}\sqrt{1+4x^2+4y^2}\,dA.
\]

In polar coordinates ($x=r\cos\theta$, $y=r\sin\theta$):
\[
\boxed{\text{Surface Area} = \int_0^{2\pi}\int_0^{\sqrt{5}}\sqrt{1+4r^2}\;r\,dr\,d\theta.}
\]

%%% PROBLEM 58 %%%
\item \textbf{Solution.}
$z=f(x,y)=x^2+4y^2$, so $f_x=2x$, $f_y=8y$.
\[
dS = \sqrt{1+f_x^2+f_y^2}\,dA = \sqrt{1+4x^2+64y^2}\,dA.
\]

On the surface, $T(x,y,z) = \sqrt{1+x^2+y^2}$ (since $z=x^2+4y^2$, $T$ depends only on $x,y$).

\[
\boxed{\iint_S T\,dS = \int_{-1}^{1}\int_{-1}^{1}\sqrt{1+x^2+y^2}\;\sqrt{1+4x^2+64y^2}\;dx\,dy.}
\]

%%% PROBLEM 59 %%%
\item \textbf{Solution.}
\begin{enumerate}
\item The plane $z=2x$ over $\{0\le x\le 2,\;0\le y\le 1\}$:
\[
\mathbf{r}(x,y)=(x, y, 2x), \quad 0\le x\le 2,\;0\le y\le 1.
\]

\item $\mathbf{r}_x=(1,0,2)$, $\mathbf{r}_y=(0,1,0)$.
\[
\mathbf{r}_x\times\mathbf{r}_y = \begin{vmatrix}\mathbf{i}&\mathbf{j}&\mathbf{k}\\ 1&0&2\\ 0&1&0\end{vmatrix} = (-2, 0, 1).
\]
For upward orientation ($z$-component positive), $\mathbf{n}\,dA = (-2,0,1)\,dx\,dy$.

$\mathbf{F}=(x^2, 0, yz)$. On the surface $z=2x$: $\mathbf{F}=(x^2, 0, 2xy)$.

\[
\boxed{\iint_S\mathbf{F}\cdot d\mathbf{S} = \int_0^1\int_0^2\left[x^2(-2)+0+2xy(1)\right]dx\,dy = \int_0^1\int_0^2(2xy-2x^2)\,dx\,dy.}
\]
\end{enumerate}

%%% PROBLEM 60 %%%
\item \textbf{Solution.}
The cylinder has radius $2$, $0\le z\le 3$.

$\mathbf{r}(\theta,z)=(2\cos\theta, 2\sin\theta, z)$.

$\mathbf{r}_\theta=(-2\sin\theta, 2\cos\theta, 0)$, $\mathbf{r}_z=(0,0,1)$.
\[
\mathbf{r}_\theta\times\mathbf{r}_z = (2\cos\theta, 2\sin\theta, 0).
\]
\[
\|\mathbf{r}_\theta\times\mathbf{r}_z\| = 2.
\]

On the cylinder, $f=\sqrt{x^2+y^2}=2$.
\[
\iint_S f\,dS = \int_0^{2\pi}\int_0^3 2\cdot 2\,dz\,d\theta = \int_0^{2\pi}\int_0^3 4\,dz\,d\theta = 4\cdot 3\cdot 2\pi = \boxed{24\pi}.
\]

%%% PROBLEM 61 %%%
\item \textbf{Solution.}
Hemisphere: $x^2+y^2+z^2=9$, $z\ge 0$, radius $R=3$.

Parameterize using spherical coordinates: $\mathbf{r}(\phi,\theta)=(3\sin\phi\cos\theta,\;3\sin\phi\sin\theta,\;3\cos\phi)$, with $0\le\phi\le\pi/2$, $0\le\theta\le 2\pi$.

$\|\mathbf{r}_\phi\times\mathbf{r}_\theta\| = 9\sin\phi$.

On the surface: $\rho = z = 3\cos\phi$.

\begin{enumerate}
\item \textbf{Total mass:}
\[
\boxed{\iint_S \rho\,dS = \int_0^{2\pi}\int_0^{\pi/2}(3\cos\phi)(9\sin\phi)\,d\phi\,d\theta = 27\int_0^{2\pi}\int_0^{\pi/2}\cos\phi\sin\phi\,d\phi\,d\theta.}
\]
Evaluating: $\int_0^{\pi/2}\cos\phi\sin\phi\,d\phi = \frac{1}{2}$, so the mass is $27\cdot2\pi\cdot\frac{1}{2}=27\pi$.

\item \textbf{Outward flux of $\mathbf{F}=(0,0,z)$:}

The outward unit normal on the sphere is $\hat{\mathbf{n}}=\frac{1}{3}(x,y,z)=(\sin\phi\cos\theta,\sin\phi\sin\theta,\cos\phi)$.

$\mathbf{F}\cdot\hat{\mathbf{n}} = z\cos\phi = 3\cos\phi\cdot\cos\phi = 3\cos^2\phi$.

\[
\boxed{\iint_S\mathbf{F}\cdot d\mathbf{S} = \int_0^{2\pi}\int_0^{\pi/2}3\cos^2\phi\cdot 9\sin\phi\,d\phi\,d\theta = 27\int_0^{2\pi}\int_0^{\pi/2}\cos^2\phi\sin\phi\,d\phi\,d\theta.}
\]
Evaluating: $\int_0^{\pi/2}\cos^2\phi\sin\phi\,d\phi = \frac{1}{3}$, so the flux is $27\cdot2\pi\cdot\frac{1}{3}=18\pi$.
\end{enumerate}

%%% PROBLEM 62 %%%
\item \textbf{Solution.}
\begin{enumerate}
\item $g(x,y)=x+2y$, so $g_x=1$, $g_y=2$.

The upward-oriented normal vector is $\mathbf{n}=(-g_x,-g_y,1)=(-1,-2,1)$.

\item $\mathbf{F}=(2x,3y,-z)$. On the surface $z=x+2y$: $\mathbf{F}=(2x, 3y, -(x+2y))$.

\begin{align*}
\mathbf{F}\cdot\mathbf{n} &= 2x(-1)+3y(-2)+(-(x+2y))(1) \\
&= -2x-6y-x-2y = -3x-8y.
\end{align*}

\[
\boxed{\iint_S\mathbf{F}\cdot d\mathbf{S} = \int_0^2\int_0^1(-3x-8y)\,dx\,dy.}
\]
\end{enumerate}

%%% PROBLEM 63 %%%
\item \textbf{Solution.}
$\mathbf{v}=(x,y,-1)$, closed surface: cylinder $x^2+y^2=4$, $-1\le z\le 2$, plus top and bottom disks.

\begin{enumerate}
\item \textbf{Curved side:} $\mathbf{r}(\theta,z)=(2\cos\theta, 2\sin\theta, z)$, $0\le\theta\le 2\pi$, $-1\le z\le 2$.

$\mathbf{r}_\theta\times\mathbf{r}_z = (2\cos\theta, 2\sin\theta, 0)$ (outward-pointing).

\item \textbf{Flux through curved side:}

$\mathbf{v}=(2\cos\theta, 2\sin\theta, -1)$ on the surface.
\[
\mathbf{v}\cdot(\mathbf{r}_\theta\times\mathbf{r}_z) = 4\cos^2\theta+4\sin^2\theta = 4.
\]
\[
\Phi_{\text{side}} = \int_0^{2\pi}\int_{-1}^{2}4\,dz\,d\theta = 4\cdot 3\cdot 2\pi = 24\pi.
\]

\textbf{Flux through top disk} ($z=2$, outward normal $\hat{\mathbf{n}}=(0,0,1)$):
\[
\mathbf{v}\cdot\hat{\mathbf{n}} = -1.
\]
\[
\Phi_{\text{top}} = \iint_{x^2+y^2\le 4}(-1)\,dA = -4\pi.
\]

\textbf{Flux through bottom disk} ($z=-1$, outward normal $\hat{\mathbf{n}}=(0,0,-1)$):
\[
\mathbf{v}\cdot\hat{\mathbf{n}} = -(-1) = 1.
\]
\[
\Phi_{\text{bot}} = \iint_{x^2+y^2\le 4}1\,dA = 4\pi.
\]

\item \textbf{Total outward flux:}
\[
\Phi = \Phi_{\text{side}}+\Phi_{\text{top}}+\Phi_{\text{bot}} = 24\pi - 4\pi + 4\pi = \boxed{24\pi}.
\]

\textbf{Verification via Divergence Theorem:} $\nabla\cdot\mathbf{v}=1+1+0=2$, and the volume is $\pi(2)^2(3)=12\pi$. So $\iiint_V 2\,dV = 24\pi$.
\end{enumerate}

%%% PROBLEM 64 %%%
\item \textbf{Solution.}
$\mathbf{F}=(y^2, z^2, x^2)$.

\textbf{Step 1:} Compute $\nabla\times\mathbf{F}$:
\[
\nabla\times\mathbf{F} = \begin{vmatrix}\mathbf{i}&\mathbf{j}&\mathbf{k}\\ \partial_x&\partial_y&\partial_z\\ y^2&z^2&x^2\end{vmatrix} = (0-2z)\mathbf{i}-(2x-0)\mathbf{j}+(0-2y)\mathbf{k} = (-2z,-2x,-2y).
\]

\textbf{Step 2:} Parameterize $\Sigma$: the paraboloid $z=x^2+y^2$ for $x^2+y^2\le 4$. Use $\mathbf{r}(x,y)=(x,y,x^2+y^2)$.

$\mathbf{r}_x=(1,0,2x)$, $\mathbf{r}_y=(0,1,2y)$.

For Stokes' Theorem with the boundary oriented counterclockwise when viewed from above, we need the upward-oriented normal:
\[
\mathbf{r}_x\times\mathbf{r}_y = (-2x,-2y,1).
\]

\textbf{Step 3:} By Stokes' Theorem:
\[
\oint_C\mathbf{F}\cdot d\mathbf{r} = \iint_\Sigma(\nabla\times\mathbf{F})\cdot d\mathbf{S}.
\]

$(\nabla\times\mathbf{F})\cdot(\mathbf{r}_x\times\mathbf{r}_y) = (-2z)(-2x)+(-2x)(-2y)+(-2y)(1)$.

On $\Sigma$: $z=x^2+y^2$:
\[
= 4x(x^2+y^2)+4xy-2y.
\]

\[
\boxed{\oint_C\mathbf{F}\cdot d\mathbf{r} = \iint_{x^2+y^2\le 4}\left[4x(x^2+y^2)+4xy-2y\right]dA.}
\]

%%% PROBLEM 65 %%%
\item \textbf{Solution.}
$\mathbf{F}=(e^y, \ln(1+z^2), y^2z)$, $\Sigma$: $x+2y+3z=6$ in the first octant.

\textbf{Decision:} It is easier to use the \textbf{surface integral} $\iint_\Sigma(\nabla\times\mathbf{F})\cdot d\mathbf{S}$ (Stokes' Theorem) rather than the line integral directly. The reason: computing $\oint_C \mathbf{F}\cdot d\mathbf{r}$ directly requires parameterizing three edges of the triangle and evaluating three separate integrals with complicated integrands ($e^y$, $\ln(1+z^2)$). In contrast, the curl may simplify, and the surface integral over a flat triangular patch is straightforward.

\textbf{Outline of steps:}
\begin{enumerate}
\item Compute $\nabla\times\mathbf{F}$:
Computing component by component:
\[
(\nabla\times\mathbf{F})_x = \frac{\partial(y^2z)}{\partial y}-\frac{\partial(\ln(1+z^2))}{\partial z} = 2yz - \frac{2z}{1+z^2},
\]
\[
(\nabla\times\mathbf{F})_y = \frac{\partial(e^y)}{\partial z}-\frac{\partial(y^2z)}{\partial x} = 0-0 = 0,
\]
\[
(\nabla\times\mathbf{F})_z = \frac{\partial(\ln(1+z^2))}{\partial x}-\frac{\partial(e^y)}{\partial y} = 0-e^y = -e^y.
\]

\item The plane $x+2y+3z=6$ has upward normal $\mathbf{n}=(1,2,3)/\sqrt{14}$ and $d\mathbf{S}=(1,2,3)\,dA'$ where $dA'$ is the projected area element.

\item Parameterize using $(y,z)$ or project onto a coordinate plane and integrate over the triangular region.
\end{enumerate}

%%% PROBLEM 66 %%%
\item \textbf{Solution.}
$\mathbf{F}=(-y,x,0)$ on the cylinder $x^2+y^2=9$, $-2\le z\le 5$.

\begin{enumerate}
\item
\[
\nabla\times\mathbf{F} = \begin{vmatrix}\mathbf{i}&\mathbf{j}&\mathbf{k}\\ \partial_x&\partial_y&\partial_z\\ -y&x&0\end{vmatrix} = (0-0)\mathbf{i}-(0-0)\mathbf{j}+(1-(-1))\mathbf{k} = (0,0,2).
\]

\item The boundary $\partial\Sigma$ consists of two circles:

\textbf{Top circle} ($z=5$): $\mathbf{r}_1(\theta)=(3\cos\theta, 3\sin\theta, 5)$, $\theta\in[0,2\pi]$, oriented counterclockwise when viewed from outside (by the right-hand rule with the outward cylinder normal).

\textbf{Bottom circle} ($z=-2$): $\mathbf{r}_2(\theta)=(3\cos\theta, 3\sin\theta, -2)$, $\theta\in[0,2\pi]$, oriented clockwise when viewed from above (i.e., counterclockwise when viewed from below, consistent with the outward-pointing normal and right-hand rule).

\item By Stokes' Theorem:
\[
\oint_{\partial\Sigma}\mathbf{F}\cdot d\mathbf{r} = \iint_\Sigma(\nabla\times\mathbf{F})\cdot d\mathbf{S}.
\]

The cylinder surface has outward normal $\hat{\mathbf{n}}=(\cos\theta, \sin\theta, 0)$.

$(\nabla\times\mathbf{F})\cdot\hat{\mathbf{n}} = (0,0,2)\cdot(\cos\theta,\sin\theta,0) = 0$.

\[
\boxed{\iint_\Sigma(\nabla\times\mathbf{F})\cdot d\mathbf{S} = 0.}
\]
\end{enumerate}

%%% PROBLEM 67 %%%
\item \textbf{Solution.}
If $\mathbf{F}=\nabla f$, then $\nabla\times\mathbf{F}=\nabla\times(\nabla f)=\mathbf{0}$ (by the identity proven in Problem 50).

By Stokes' Theorem:
\[
\oint_C \mathbf{F}\cdot d\mathbf{r} = \iint_\Sigma(\nabla\times\mathbf{F})\cdot d\mathbf{S} = \iint_\Sigma\mathbf{0}\cdot d\mathbf{S} = 0.
\]

Therefore, for any smooth oriented surface $\Sigma$ with boundary $C$:
\[
\oint_C \nabla f\cdot d\mathbf{r} = 0. \qquad\blacksquare
\]

This is consistent with the fact that conservative fields have path-independent line integrals, so the integral around any closed loop is zero.

%%% PROBLEM 68 %%%
\item \textbf{Solution.}
Faraday's Law states:
\[
\oint_C \mathbf{E}\cdot d\mathbf{r} = -\frac{d}{dt}\iint_\Sigma\mathbf{B}\cdot d\mathbf{S}.
\]

By Stokes' Theorem, the left side equals:
\[
\oint_C \mathbf{E}\cdot d\mathbf{r} = \iint_\Sigma(\nabla\times\mathbf{E})\cdot d\mathbf{S}.
\]

If the surface $\Sigma$ is fixed (static boundary), we can bring the time derivative inside:
\[
\iint_\Sigma(\nabla\times\mathbf{E})\cdot d\mathbf{S} = -\iint_\Sigma\frac{\partial\mathbf{B}}{\partial t}\cdot d\mathbf{S}.
\]

Since this holds for every surface $\Sigma$, the integrands must be equal:
\[
\nabla\times\mathbf{E} = -\frac{\partial\mathbf{B}}{\partial t},
\]
which is the differential form of Faraday's Law.

Here $\mathbf{E}$ plays the role of the vector field in the line integral (circulation of the electric field), and $\nabla\times\mathbf{E}$ plays the role of the flux integrand (the curl relates to the surface integral via Stokes).

\textbf{Physical meaning of the negative time derivative:} A \emph{changing} magnetic flux through $\Sigma$ induces an electric field that circulates around $C$. The negative sign (Lenz's law) means the induced electric field opposes the change in magnetic flux: if $\mathbf{B}$ increases, the induced $\mathbf{E}$ circulates in the direction that would create a magnetic field opposing the increase.

%%% PROBLEM 69 %%%
\item \textbf{Solution.}
$\mathbf{F}=(xy, yz, xz)$.
\[
\nabla\cdot\mathbf{F} = \frac{\partial(xy)}{\partial x}+\frac{\partial(yz)}{\partial y}+\frac{\partial(xz)}{\partial z} = y+z+x.
\]

By the Divergence Theorem:
\[
\oiint_{\partial V}\mathbf{F}\cdot\mathbf{n}\,dS = \iiint_V(x+y+z)\,dV,
\]
where $V$ is the ball $x^2+y^2+z^2\le 4$.

In spherical coordinates: $x+y+z = \rho\sin\phi\cos\theta+\rho\sin\phi\sin\theta+\rho\cos\phi$, $dV=\rho^2\sin\phi\,d\rho\,d\phi\,d\theta$.

By symmetry, $\iiint_V x\,dV = \iiint_V y\,dV = \iiint_V z\,dV = 0$ (each is an odd function integrated over a symmetric domain).

\[
\boxed{\iiint_V(x+y+z)\,dV = 0.}
\]

%%% PROBLEM 70 %%%
\item \textbf{Solution.}
$\mathbf{F}=(x^2,y^2,z^2)$. The region $V$: $0\le z\le\sqrt{x^2+y^2}$ capped by $z=2$.

More precisely, $V$ is the region inside the cone $z=\sqrt{x^2+y^2}$ and below $z=2$. In cylindrical coordinates: $r\le z\le 2$, $0\le r\le 2$, $0\le\theta\le 2\pi$.

\begin{enumerate}
\item \textbf{Boundary surfaces:}
\begin{itemize}
\item The cone: $z=\sqrt{x^2+y^2}$ (i.e., $z=r$) for $0\le r\le 2$ (with inward-pointing normal to close the volume, i.e., normal pointing outward from $V$, which is downward-and-outward on the cone).
\item The top disk: $z=2$, $x^2+y^2\le 4$ (outward normal is $\hat{\mathbf{n}}=(0,0,1)$).
\end{itemize}

\item By the Divergence Theorem:
\[
\nabla\cdot\mathbf{F} = 2x+2y+2z.
\]
\[
\iint_{\partial V}\mathbf{F}\cdot d\mathbf{S} = \iiint_V(2x+2y+2z)\,dV.
\]

In cylindrical coordinates:
\[
= \int_0^{2\pi}\int_0^2\int_r^2(2r\cos\theta+2r\sin\theta+2z)\,r\,dz\,dr\,d\theta.
\]

By symmetry, the $\cos\theta$ and $\sin\theta$ terms integrate to zero over $[0,2\pi]$:
\[
\boxed{= \int_0^{2\pi}\int_0^2\int_r^2 2zr\,dz\,dr\,d\theta = 2\pi\int_0^2 r\left[z^2\right]_r^2 dr = 2\pi\int_0^2 r(4-r^2)\,dr.}
\]
Evaluating: $\int_0^2 r(4-r^2)\,dr = \int_0^2(4r-r^3)\,dr = [2r^2-r^4/4]_0^2 = 8-4=4$. So the total flux is $8\pi$.
\end{enumerate}

%%% PROBLEM 71 %%%
\item \textbf{Solution.}
$\mathbf{F}=(e^z, \sin(xy), 2z)$, $V$: cylinder $x^2+y^2\le 1$, $0\le z\le 5$.

\textbf{Decision:} Use the \textbf{Divergence Theorem} (volume integral), because computing the surface integral directly requires parameterizing three surfaces (curved side, top disk, bottom disk) with complicated integrands like $e^z$ and $\sin(xy)$.

\textbf{Divergence:}
\[
\nabla\cdot\mathbf{F} = 0 + x\cos(xy) + 2.
\]

By the Divergence Theorem:
\[
\oiint_{\partial V}\mathbf{F}\cdot\mathbf{n}\,dS = \iiint_V(x\cos(xy)+2)\,dV.
\]

\textbf{Steps:}
\begin{enumerate}
\item The integral of $2$ over $V$: $2\cdot\text{Vol}(V) = 2\cdot\pi(1)^2(5)=10\pi$.

\item The integral of $x\cos(xy)$ over $V$: Note $x\cos(xy)$ is an \emph{odd} function of $x$ (replacing $x\to -x$ gives $-x\cos(-xy)=-x\cos(xy)$), and the domain is symmetric in $x$. Therefore $\iiint_V x\cos(xy)\,dV = 0$.
\end{enumerate}

\[
\boxed{\oiint_{\partial V}\mathbf{F}\cdot\mathbf{n}\,dS = 10\pi.}
\]

%%% PROBLEM 72 %%%
\item \textbf{Solution.}
$\mathbf{F}=(x^2+y^2, z, 3x)$, $V$: ball $x^2+y^2+z^2\le 4$.
\[
\nabla\cdot\mathbf{F} = 2x+0+0 = 2x.
\]

In spherical coordinates: $x = \rho\sin\phi\cos\theta$, $dV=\rho^2\sin\phi\,d\rho\,d\phi\,d\theta$.
\[
\iiint_V 2x\,dV = \int_0^{2\pi}\int_0^{\pi}\int_0^{2} 2\rho\sin\phi\cos\theta\cdot\rho^2\sin\phi\,d\rho\,d\phi\,d\theta.
\]
\[
\boxed{= 2\int_0^{2\pi}\cos\theta\,d\theta\int_0^{\pi}\sin^2\phi\,d\phi\int_0^2\rho^3\,d\rho.}
\]

Since $\int_0^{2\pi}\cos\theta\,d\theta = 0$, the entire integral equals $\boxed{0}$.

%%% PROBLEM 73 %%%
\item \textbf{Solution.}
If a fluid is \textbf{incompressible}, then $\nabla\cdot\mathbf{v}=0$ everywhere. This means at every point, the rate of fluid creation (or destruction) per unit volume is zero---fluid is neither created nor destroyed anywhere.

By the Divergence Theorem, for any closed surface $S=\partial V$ bounding a volume $V$:
\[
\oiint_{\partial V}\mathbf{v}\cdot d\mathbf{S} = \iiint_V\nabla\cdot\mathbf{v}\,dV = \iiint_V 0\,dV = 0.
\]

The net outward flux through any closed surface is zero. Physically, this means whatever amount of fluid enters the volume through $S$ must also leave through $S$ (conservation of mass with constant density). No fluid accumulates inside, and none disappears. This is the mathematical statement of incompressibility: the volume of any fluid parcel is preserved over time.

\end{enumerate}
